\section{Artifact and evaluation}

\subsection{Artifact scope}

We evaluate the clean reference repository rather than the larger framework effort that motivated it. The repository is intentionally small. It contains the compile-time policy model, the Scala~3 macro derivation for \texttt{SchemaConforms}, Spark schema derivation from contract types, a policy-aware runtime comparator, a typed builder path that enforces sink checks, and a reproducible benchmark harness. This makes the evaluation a mechanism evaluation: the question is whether the repository proves the structural properties it claims, not whether it already delivers industrial-scale outcomes.

\subsection{Compile-time and runtime coverage}

Compile-time coverage is centered on direct witness generation and compile-fail behavior. The tests exercise positive conformance for exact-style, backward-compatible, and forward-compatible policies, and they exercise rejection paths for reordering, missing required fields, and nested optionality drift. The shape family covered by those tests includes nested case classes, sequences, maps, optional fields, and defaulted contract fields.

The runtime layer is evaluated separately because it has a different job. The compile-time witness proves compatibility for declared Scala types. The runtime layer validates the actual Spark schema that reaches the sink. The runtime tests therefore focus on three questions: does the exact-style pin reject nested collection optionality drift that Spark ignores by default, do \texttt{Backward} and \texttt{Forward} behave as true subset relations at runtime, and does runtime checking use contract-derived metadata to allow missing defaulted fields where the policy permits them.

\subsection{Builder-path end-to-end checks}

The builder-path tests are the most operational part of the evaluation. They show that the paper is not only about a macro witness in isolation. \texttt{PipelineBuilder.addSink[R, P]} requires compile-time evidence for the current producer type, and the built pipeline still validates the actual \texttt{DataFrame} schema before write. The end-to-end checks cover a compile-time rejection at the sink, a green path where both the declared type and runtime schema satisfy the sink policy, a red path where the runtime schema violates the policy after compile time has succeeded, and no-transform subset-policy paths for \texttt{Backward} and \texttt{Forward}.

\subsection{Benchmark harness}

The benchmark harness is intentionally small. It measures compile-time witness-generation overhead on synthetic but representative schemas and runtime comparator overhead on nested \texttt{StructType} comparisons. It does not claim end-to-end Spark job performance.

The saved snapshots show compile-time deltas of \texttt{+0.270s}, \texttt{+0.397s}, and \texttt{+0.513s} on local \texttt{macOS arm64} for \texttt{10}, \texttt{25}, and \texttt{50} schema pairs. On GitHub-hosted \texttt{Ubuntu x86\_64}, the same sizes measure \texttt{+0.847s}, \texttt{+1.000s}, and \texttt{+1.880s}. Runtime exact unordered matching stays in the low-microsecond range per schema comparison on both environments, while exact-by-position matching stays below microsecond scale in both saved runs.

The right interpretation is limited but useful: the harness reproduces across two environments and shows that the mechanism is measurable and stable enough to package as artifact evidence. The wrong interpretation would be a broad claim about low overhead across machines or organizations. Table~\ref{tab:benchmarks} therefore reports the saved snapshots as bounded artifact evidence, not as a general performance claim.

\begin{table}[t]
  \caption{Saved benchmark snapshots for the artifact head. These numbers show reproducible harness output, not a broad cross-machine performance baseline.}
  \Description{Two-part benchmark table. The left half reports compile-time deltas for 10, 25, and 50 schema pairs on local macOS arm64 and GitHub-hosted Ubuntu x86_64. The right half reports runtime comparator averages for by-position matching, unordered exact matching, and two Spark baseline comparators on the same two environments.}
  \label{tab:benchmarks}
  \centering
  \scriptsize
  \textbf{Compile-time overhead}

  \begin{tabular}{lrrrr}
    \toprule
    Pairs & Local (s) & Local (\%) & Ubuntu (s) & Ubuntu (\%) \\
    \midrule
    10 & 0.270 & 11.8 & 0.847 & 12.6 \\
    25 & 0.397 & 13.5 & 1.000 & 11.1 \\
    50 & 0.513 & 13.9 & 1.880 & 16.6 \\
    \bottomrule
  \end{tabular}

  \vspace{0.6em}
  \textbf{Runtime comparator averages}

  \begin{tabular}{lrr}
    \toprule
    Benchmark & Local & Ubuntu \\
    \midrule
    By-position & 116.82 & 180.55 \\
    Unordered exact & 4736.41 & 8149.74 \\
    Spark ignore-case & 278.92 & 331.42 \\
    Spark structural & 332.13 & 380.36 \\
    \bottomrule
  \end{tabular}
\end{table}
